\documentclass[11pt]{article}

\usepackage{arxivrefined}
\usepackage{amsmath,amssymb,booktabs,array}
\usepackage[hidelinks]{hyperref}
\usepackage{enumitem}
\usepackage{microtype}

\newcommand{\method}{FI-JEPA}

\title{FI-JEPA on a Compact Macroeconomic Forecasting Task:\\An Evidence-Bounded Negative Result for Learned Latent Prediction}
\author{Ryan Gomez}
\date{September 2026}

\begin{document}
\maketitle

\begin{abstract}
We study FI-JEPA, a joint-embedding predictive architecture intended to learn latent financial dynamics rather than directly forecast raw observations. The implementation combines an online encoder, an exponential-moving-average target encoder, optional prototype memory, stochastic latent prediction, and finance-motivated regularization. We evaluate a frozen, bounded real-data protocol on the public \texttt{statsmodels} macrodata panel using three predeclared training seeds, chronological train/validation/test splits, validation-only checkpoint selection, and fixed downstream Ridge probes. The retained result is negative for the maintained representation claim on this task: a deterministic Ridge probe on flattened raw context substantially outperforms the same probe applied to the learned FI-JEPA latent representation, and removing several proposed components does not degrade mean regression error. A post-run audit also finds that the planned split-versus-monolithic predictor control is non-identifying under the frozen one-stage macro configuration, so it is retained for provenance but excluded from mechanism inference. We report the adverse result without rescue tuning. The study is deliberately narrow: it does not establish a general conclusion about JEPA methods, market prediction, trading performance, or financial representation learning.
\end{abstract}

\section{Introduction}
Financial time series are noisy, nonstationary, and influenced by mechanisms operating on different temporal scales. This motivates representation-learning approaches that attempt to model persistent latent structure rather than directly fit every fluctuation in observation space. FI-JEPA was designed around that intuition. A context encoder maps a recent market window into a latent state; a slowly moving target encoder defines future latent targets; a predictor advances the latent state; and auxiliary losses attempt to encode financially motivated structure.

That motivation alone does not establish that the representation is useful. A learned latent representation should be compared against simple controls on a frozen downstream task, and architecture components should only receive mechanistic credit when their removal changes executed computation and produces reproducible evidence. This paper therefore treats the compact macrodata experiment as a falsifiable representation-learning study rather than as a demonstration exercise.

The present protocol was frozen before its canonical multi-seed run. The result is adverse to the central bounded claim. On the retained GDP-growth proxy task, the raw-context linear control is substantially better than the FI-JEPA latent probe. Several component removals also have equal or lower mean regression error than the full model. In addition, a post-run code audit establishes that one planned operator-factorization control was structurally non-identifying because both arms instantiated a single predictor stage. We preserve that failed control rather than silently replacing it after observing outcomes.

The contributions of this report are therefore primarily methodological and negative-result oriented:
\begin{enumerate}[leftmargin=1.5em]
    \item a reproducible real-data FI-JEPA training and evaluation path with validation-selected checkpoints and per-seed provenance;
    \item a fixed raw-context Ridge downstream control using the same target, split, and probe family as the learned-latent evaluation;
    \item real component removals for EMA updating, finance-aware regularization, stochastic uncertainty heads, and prototype memory;
    \item an explicit post-run identification of a non-identifying operator-split control;
    \item an evidence-bound preprint in which the result table is generated directly from the canonical retained JSON artifact.
\end{enumerate}

\section{Method}
\subsection{Context and target encoders}
Let $X_t \in \mathbb{R}^{L\times d}$ denote a context window of length $L$ with $d$ observed features. An online encoder $f_\theta$ maps the window to a latent state
\begin{equation}
    z_t = f_\theta(X_t) \in \mathbb{R}^{m}.
\end{equation}
A target encoder $f_{\bar\theta}$ has the same architecture but does not receive ordinary gradient updates. In the full condition its parameters follow an exponential moving average,
\begin{equation}
    \bar\theta \leftarrow \tau \bar\theta + (1-\tau)\theta.
\end{equation}
The \texttt{no\_ema} condition instead copies the online encoder state into the target path after optimizer updates. This changes the executed target-update mechanism while retaining the rest of the training path.

\subsection{Latent predictor}
A predictor stage receives a latent state, transforms it through a small multilayer perceptron, and outputs a latent mean $\mu$, a gate $g$, and, when stochastic prediction is enabled, a log variance $\log\sigma^2$. A stochastic stage samples
\begin{equation}
    z' = \mu + g \odot \sigma \odot \epsilon, \qquad \epsilon \sim \mathcal{N}(0,I),
\end{equation}
while the deterministic condition uses the mean transition without the stochastic head. The canonical macro configuration enables stochastic prediction, so \texttt{no\_uncertainty\_heads} is an executed intervention rather than a configuration alias.

The general FI-JEPA implementation supports a sequence of predictor stages. However, the frozen macro protocol used \texttt{stage\_names: [macro]}. Consequently, its nominal split predictor contains only one stage. Switching to a nominal monolithic predictor also creates one stage with the same numerical architecture. The stage name itself is not consumed by the forward computation. The resulting split-versus-monolithic row is therefore not a valid operator-factorization ablation in this protocol. This fact was identified after the canonical run and is not repaired post hoc.

\subsection{Prototype memory}
The optional memory stores prototype keys and values and retrieves a weighted context for the current latent state. A learned gate mixes the retrieved memory into the context representation. The \texttt{no\_memory} condition removes this memory path entirely. Memory-update errors are fail-closed in the audited runner rather than silently swallowed.

\subsection{Training objective}
The implementation combines latent prediction with optional uncertainty and finance-aware terms. In schematic form,
\begin{equation}
\mathcal{L} = \mathcal{L}_{\mathrm{pred}}
+ \lambda_{\mathrm{nll}}\mathcal{L}_{\mathrm{nll}}
+ \lambda_{\mathrm{fin}}\mathcal{L}_{\mathrm{fin}}
+ \lambda_{\mathrm{aux}}\mathcal{L}_{\mathrm{aux}}.
\end{equation}
The \texttt{no\_financial\_regularizers} condition disables the finance-aware regularization group while preserving the latent-prediction path. We do not interpret a component as beneficial merely because it has a financially suggestive name; its value must appear in the frozen comparison.

\section{Frozen Real-Data Protocol}
The canonical protocol identifier is \texttt{FIJEPA\_MACRODATA\_PAPER\_V1\_20260909}. It is documented in \texttt{research/PAPER\_PROTOCOL\_V1.md}, and the retained machine-readable output is \texttt{experiments/paper\_results.json}.

\subsection{Dataset and target}
The study uses the public quarterly macroeconomic panel distributed as \texttt{statsmodels.datasets.macrodata}. The target is a GDP-growth proxy formed from real GDP. The context length is four quarters and the forecast horizon is one quarter. Data are split chronologically into 70\% train, 15\% validation, and 15\% test partitions. Normalization statistics are fit on the training partition only.

Window construction occurs independently inside each partition. This is conservative at split boundaries because it discards cross-boundary history instead of allowing later-partition information to enter an earlier training example.

\subsection{Seeds, optimization, and model selection}
The predeclared seeds are 7, 17, and 27. Every model condition receives the same 15-epoch budget, batch-size policy, optimizer family, learning rate, data split, and downstream probe implementation. Each variant and seed has an isolated run directory so histories and checkpoints cannot overwrite one another.

The paper-facing representation is loaded from the checkpoint with minimum validation total loss. The test partition is not used for checkpoint selection. Non-finite losses, non-finite gradient norms, empty loaders, and memory-update failures are treated as errors rather than converted into successful runs.

\subsection{Conditions}
The retained matrix contains the full model and the following planned conditions:
\begin{itemize}[leftmargin=1.5em]
    \item \texttt{no\_ema}: remove the moving-average target update;
    \item \texttt{no\_financial\_regularizers}: remove finance-aware regularizers;
    \item \texttt{no\_uncertainty\_heads}: remove stochastic latent uncertainty;
    \item \texttt{no\_memory}: remove prototype memory;
    \item \texttt{no\_operator\_split}: nominal monolithic predictor control.
\end{itemize}
The first four are executed component removals. The last is retained only as a non-identifying control because the frozen full configuration contains a single predictor stage.

A historical \texttt{deterministic} alias is not counted as a separate ablation because it is implementation-equivalent to \texttt{no\_uncertainty\_heads}.

\subsection{Matched downstream control}
For each seed, FI-JEPA embeddings are evaluated with \texttt{Ridge(alpha=1.0)} for the continuous GDP-growth target and the corresponding fixed Ridge classifier for a median-thresholded target. The raw-context control applies the same downstream estimators to flattened context windows from the same chronological rows and target construction. This isolates whether the learned embedding improves the fixed probe relative to directly exposing the probe to the recent observations.

The raw-context Ridge is a matched downstream-evaluation control. It is not a parameter-matched neural architecture baseline and is not presented as one.

\section{Results}
All numerical entries in Table~\ref{tab:canonical_macro} are generated from the retained canonical JSON artifact by \texttt{scripts/make\_canonical\_paper\_assets.py}. Historical one-seed tables and figures are not evidence sources for this preprint.

% AUTO-GENERATED from experiments/paper_results.json; do not hand edit.
\begin{table}[t]
\centering
\caption{Canonical three-seed macrodata study. Mean $\pm$ sample standard deviation across predeclared seeds. Regression MSE/MAE are lower-is-better; directional/classification accuracy are higher-is-better. The monolithic row is retained for provenance but is non-identifying because the frozen full macro configuration contains only one executed predictor stage.}
\label{tab:canonical_macro}
\begin{tabular}{lcccc}
\toprule
Condition & Regression MSE & Regression MAE & Directional acc. & Classification acc. \\
\midrule
Full FI-JEPA & 0.920145 $\pm$ 0.074493 & 0.683738 $\pm$ 0.048334 & 0.259259 $\pm$ 0.064150 & 0.234568 $\pm$ 0.021383 \\
No EMA & 0.872222 $\pm$ 0.033997 & 0.650183 $\pm$ 0.032533 & 0.259259 $\pm$ 0.064150 & 0.234568 $\pm$ 0.021383 \\
No financial regularizers & 0.915821 $\pm$ 0.071222 & 0.681149 $\pm$ 0.045780 & 0.259259 $\pm$ 0.064150 & 0.234568 $\pm$ 0.021383 \\
Monolithic predictor (non-identifying control) & 0.920145 $\pm$ 0.074493 & 0.683738 $\pm$ 0.048334 & 0.259259 $\pm$ 0.064150 & 0.234568 $\pm$ 0.021383 \\
No uncertainty heads & 0.912092 $\pm$ 0.051924 & 0.682212 $\pm$ 0.043498 & 0.271605 $\pm$ 0.085533 & 0.234568 $\pm$ 0.021383 \\
No memory & 0.895743 $\pm$ 0.044221 & 0.668672 $\pm$ 0.027983 & 0.271605 $\pm$ 0.085533 & 0.234568 $\pm$ 0.021383 \\
Raw-context Ridge & 0.441924 $\pm$ 0.000000 & 0.537767 $\pm$ 0.000000 & 0.629630 $\pm$ 0.000000 & 0.777778 $\pm$ 0.000000 \\
\bottomrule
\end{tabular}
\end{table}


The central result is negative. The learned full-model latent representation performs substantially worse than the raw-context Ridge control across regression error, directional accuracy, and the retained binary classification probe. This comparison is particularly important because the control is simple and receives no learned FI-JEPA representation.

The ablation pattern also does not provide support for the proposed components on this task. Mean regression MSE is lower after removing the EMA update, finance-aware regularizers, or prototype memory. Removing uncertainty heads is mixed across seeds and does not establish a benefit for the full stochastic model. These observations are descriptive because only three training seeds were predeclared.

The nominal monolithic-predictor row is numerically identical to the full row in the retained table. That equality is not interpreted as evidence that operator factorization is neutral. The post-run source audit shows that the two arms are implementation-equivalent under the frozen one-stage macro configuration, making the comparison non-identifying.

\section{Interpretation}
The bounded hypothesis that FI-JEPA's learned latent representation improves the fixed downstream GDP-growth probe is not supported here. The simplest explanation consistent with the retained evidence is that, on this small quarterly panel and training budget, the representation learner discards or distorts predictive information that the linear probe can exploit directly from the raw context.

The component results also caution against attributing value to architectural complexity without isolation. EMA targets, finance-aware regularization, memory, and stochastic prediction may be useful under other tasks or data regimes, but this protocol does not establish those benefits. In particular, the finance-aware losses should not be treated as validated economic structure merely because they are motivated by finance.

The adverse result is scientifically useful because it narrows the next research question. Future work should ask under what data scale, task structure, or latent objective a learned predictive representation can beat strong simple controls, rather than assuming that a more structured architecture should win. Any such successor must be separately frozen before new outcome access; it cannot retroactively rescue protocol v1.

\section{Limitations}
The most important limitations are:
\begin{enumerate}[leftmargin=1.5em]
    \item \textbf{Small dataset.} The macrodata panel is a compact quarterly dataset and cannot establish broad financial generalization.
    \item \textbf{Three seeds.} The retained paired bootstrap intervals summarize only three training seeds and are descriptive rather than population-level significance guarantees.
    \item \textbf{Narrow target.} The GDP-growth proxy is not a trading, portfolio, liquidity, or market-microstructure objective.
    \item \textbf{Baseline scope.} Raw-context Ridge is a strong simple downstream control for the representation question but is not a compute- or parameter-matched neural baseline.
    \item \textbf{Non-identifying planned control.} The operator split cannot be evaluated from protocol v1 because both frozen arms execute one predictor stage.
    \item \textbf{No external market validation.} The study provides no evidence about profitability, live trading, transaction costs, or deployment robustness.
\end{enumerate}

\section{Reproducibility and Evidence Boundary}
The canonical result was produced by a GitHub Actions workflow from the frozen branch and retained back into the repository with a run receipt. The evidence gate requires the declared seed matrix, finite metrics, an explicit raw-context control, and paired seed-level summaries. The runner uses validation-selected checkpoints and isolated output directories.

The post-run audit identified a weakness in the original structural signature check: it treated predictor mode and stage labels as sufficient to call the operator control distinct even though both configurations executed one numerical stage. That defect is now documented as a validity limitation rather than hidden. The raw v1 artifact is preserved unchanged.

This paper authorizes none of the following claims: trading alpha, profitability, market-wide generalization, state-of-the-art forecasting, universal failure of JEPA, universal failure of financial self-supervision, or causal benefit/harm of the non-identifying operator-factorization control.

\section{Conclusion}
In a frozen three-seed real-data macroeconomic study, FI-JEPA's learned latent representation does not beat a simple raw-context Ridge control, and the retained component removals do not establish benefits for several proposed mechanisms. A planned operator-factorization ablation is additionally non-identifying under the frozen one-stage configuration. We preserve these adverse findings without post-outcome rescue tuning. The appropriate scientific conclusion is narrow and negative: this FI-JEPA configuration does not demonstrate useful representation improvement on the tested compact macrodata task. Stronger claims require a separately preregistered study with an identifying architecture, larger and more varied data, and appropriately matched baselines.

\end{document}
